\section{Scope, source audit, and the contribution}
\label{sec:scope}

The purpose of this note is to make mathematical work in \ipra\ more
composable while actually improving one of the repository's interfaces.
The supplied repository is substantial: the inspected snapshot contains
sixty theorem packets, including nine earlier frameworks. Their index
explicitly distinguishes the checked ConciseIPRA development from packets
filed as delivered without a repository-wide verification pass
\cite{RepoIndex,PacketIndex}. Repeating their finite-observation vocabulary
would not resolve an additional mathematical obligation.

We therefore pin the audit to commit
\begin{center}\small
\path{b0f6f535ffc16324683fef2c263260a88ae03755}
\end{center}
and concentrate on two missing connections. The first connects a local
proof about a change of representation to a derivation about an entire
finite computation. The second connects an algebraic certificate of
finite dimensionality to positive arithmetic witnesses: coordinate bounds
and complete rational lists.

\subsection{Results, in their exact logical forms}

The main logical result has a deliberately syntactic premise. Let
$f,g,h$ be fixed unary PR function symbols. From a supplied derivation
in \ipra\ of
\[
                      h(f(x))=g(h(x))
\]
we construct a derivation of
\[
       \forall n\;h(\iter_f(n,x))=\iter_g(n,h(x)).
\]
The proof uses the displayed equality itself as a quantifier-free
induction invariant. The companion theorem transports an arbitrary
quantifier-free binary relation preserved by a pair of steps.
These are checked Lean theorems with \code{Prf PRAAx} conclusions;
they are not a rule promoting a theorem of the host logic to an
arithmetic axiom. A concrete repeated-squaring corollary connects this
infrastructure to polynomial evaluation, modular computations and
truncated series.

The main algebraic result starts with polynomial identities, all checked
by finite coefficient arithmetic. Suppose $F_1,\ldots,F_s$ are rational
polynomials in $r$ variables and we are supplied a finite family $b_i$,
a unit combination, and multiplication matrices $M_k$. At a common
zero $p$, the evaluation row $v=(b_i(p))$ is nonzero and satisfies
\[
                           vM_k=p_kv.
\]
Choose a positive integer $\Delta_k$ clearing the entries of $M_k$,
and put
\[
 A_k=\Delta_kM_k,\qquad
 R_k=\max_j\sum_i |(A_k)_{ij}|.
\]
We prove
\[
             p_k\in\{z/\Delta_k:-R_k\le z\le R_k\}.
\]
Consequently exhaustive exact filtering of at most
$\prod_k(2R_k+1)$ tuples returns \emph{all} rational common zeros.
No height bound is an extra input to this theorem: the accepted matrices
compute it. An abstract assertion that there are finitely many rational
zeros is not an adequate input.

The norm result addresses a second existing interface. For a monic plane
equation $F(X,Y)$ and a supplied section $G(X,Y)$, compute the norm of
$G$ in the free $\Q[X]$-module $\Q[X,Y]/(F)$. If that polynomial
norm is nonzero, bounded rational-root search gives a complete list of
the common rational zeros. Thus a known section can be located, not
merely counted. A global supplied finite covering list of sections can
then be enumerated by finite union. A tuple-dependent covering assertion
does not provide such a list.

\subsection{Relationship with the existing frameworks}

\begin{longtable}{P{.23\linewidth}P{.34\linewidth}P{.33\linewidth}}
\toprule
Inspected interface & Existing content & Additional result here\\
\midrule\endhead
Finite contracts and observations & Booleanized bounded invariants,
precision demands and guarded constructions. & A compiled transformation
of local object derivations into arbitrary finite simulations.\\[4pt]
Coherent methods & Finite uniqueness, exact restriction maps, filtered
lifting and object-level contract composition. & A general iteration
transport theorem and its relational version.\\[4pt]
Finite algebra certificates & A size-$D$ spanning certificate bounds
the length of a list of distinct rational zeros by $D$. & The same
action matrices give denominator divisibility, coordinate bounds and
complete rational enumeration.\\[4pt]
Finite norm certificates & Polynomial norms control the number of
points in a supplied algebraic section. & A nonzero norm gives a
complete section list by two finite rational-root searches.\\[4pt]
Finite covers in the Siegel route & Tuple-local proper-subspace covers
give a cardinality implication. & A precise stronger input, an actual
global list or a coefficient bound for the sections, gives enumeration.\\
\bottomrule
\end{longtable}

The direct finite-algebra predecessor is
\path{foundations/ipra_embedding_framework/paper/05_diophantine.tex},
especially its subsection on finite algebra certificates
\cite{Embedding}. The most recent norm and cover interfaces occur in
\path{ipra_finite_demand} and \path{ipra_coherent_methods_2026-11-15}
\cite{Demand,Coherent}. The earlier
\path{ipra_finite_observation/tex/sections.tex} explicitly explains
why cardinality and finite-tuple certificates alone yield no global
height cutoff \cite{Observation}. We use those distinctions and
strengthen the data-rich algebraic cases; we do not contradict the
general obstruction.

Multiplication matrices and effective solving of zero-dimensional
polynomial systems are established methods in computer algebra.
Rouillier's rational univariate representation is an important example
\cite{Rouillier}. This paper does not claim that these broad ideas are
new. The contribution relative to the inspected repository is the
small certificate-to-grid theorem, its weak-arithmetic proof,
the explicit constructors and scope distinctions, and the checked
object-language transport component.

\subsection{Three levels of verification}

We distinguish the following assertions throughout.
\begin{enumerate}
\item A \emph{human \ipra\ argument} describes fixed PR algorithms and
proves them using finite identities and explicit quantifier-free
induction invariants. The new algebraic and norm results have such
arguments here.
\item A \emph{host Lean theorem} proves an assertion about natural
numbers, integers, lists or rational arithmetic in Lean. The compiled
integer-eigenvalue lemmas and finite-filter reduction have this status.
\item An \emph{object Lean theorem} has a conclusion in the inductively
defined relation \code{Prf PRAAx}. The compiled simulation, fusion and
repeated-squaring transport theorems have this stronger, syntactic status.
\end{enumerate}
The full rational enumeration theorem has not yet been reified as a
\code{Prf PRAAx} derivation. Its Lean file defines the finite data and
programs and states the remaining obligations as named propositions,
without axioms or placeholder proofs. Section~\ref{sec:formalization}
gives the dependency order and the precise final arithmetic formula.

\subsection{What this does not resolve}

The reported difficulty in general Mordell--Weil is producing and
certifying the covering representatives, not running a descent after
they have been supplied \cite{MordellWeil}. The general Siegel and
Faltings packets do not supply the global effective data required for
their original requests \cite{Siegel,Faltings}. The analytic
prime-distribution packet leaves substantial uniform estimates open
\cite{PrimeDistribution}. The present framework supplies neither those
estimates nor a hidden choice procedure. It does supply new, fully
proved algebraic subroutines and makes a completed formal arithmetic
interface closer for them.

The audit is targeted. We inspected the repository index, all nine
framework descriptions, their relevant sections, the directly affected
arithmetic-geometry reports, and the actual ConciseIPRA definitions.
We did not independently verify every theorem in all sixty packets.
